\section{讨论 Discussion}\label{sec:discussion}

直接的从头算拟合改变了恒星推断中两类学习方法的角色。光谱模拟器对物理光谱库做插值，数据驱动流量模型从巡天观测中学习；二者都在恒星标签与被求值的光谱之间插入一个学习到的“标签到流量”关系。当合成以秒计、收敛大气以秒到数十秒计，物理前向模型就可以留在推断、谱线定标与验证内部。学习型大气模型可加速探索，而物理收敛仍决定被接受的大气与光谱。这改变了我们使用模拟器的方式，也改变了我们研究“观测—模型差距”的方式。

\begin{EN}
Direct ab initio fitting changes the role of two learned approaches in stellar
inference. Spectral emulators interpolate libraries of
physical spectra, while data-driven flux models learn from survey observations.
Both classes insert a learned label-to-flux relation between the stellar labels and the
spectrum being evaluated. When synthesis takes seconds and a converged
atmosphere takes seconds to tens of seconds, the physical forward model can
instead remain inside inference, line calibration, and validation. Learned
atmosphere models can accelerate exploration while physical convergence still
determines the accepted atmosphere and spectrum. This changes how we use
emulators and how we investigate the observation--model gap.
\end{EN}

\subsection{直接合成与光谱模拟器的角色 Direct synthesis and the role of spectral emulators}\label{sec:discussion_emulators}

物理自洽的丰度分析应该把相关恒星标签一起变化并拟合全光谱。一种元素改变的不只是归属它的谱线：镁改变电子供应从而改变 H$^-$ 连续谱；分子平衡与谱线覆盖在各物种间重新分配流量，并可改变大气结构 \citep{John1988,Gustafsson2008,Meszaros2012,Ting2018Oxygen,Matsuno2024}。只用某元素的具名特征测量该丰度、同时冻结模型的其余部分，会错过这些耦合响应。

\begin{EN}
A physically consistent abundance analysis should vary the relevant stellar
labels together and fit the full spectrum. An element changes more than the
lines assigned to it. Magnesium changes the electron supply and therefore the
H$^-$ continuum. Molecular balance and line blanketing redistribute flux among
species and can alter the atmospheric structure
\citep{John1988,Gustafsson2008,Meszaros2012,Ting2018Oxygen,Matsuno2024}.
Measuring each abundance only from its named features while holding the rest of
the model fixed misses these coupled responses.
\end{EN}

困难在于维度：每增加一个丰度，标签空间扩大一维，规则插值网格迅速稀疏 \citep{Ting2016}。Starfish、The Payne 及后续神经架构等概率与神经模拟器用物理光谱的平滑插值替代了网格 \citep{Czekala2015,Ting2019,Rozanski2025a}。我们当年提出 The Payne，是为了让“从自洽从头算模型做同时全光谱推断”在计算上可行——神经模拟器是加速该计算的手段，而非其科学目标。但这个名字后来成了模拟器本身的代称，而非物理拟合框架。

\begin{EN}
The difficulty is dimensionality. Each abundance expands the label space, and
regular interpolation grids become rapidly sparse \citep{Ting2016}.
Probabilistic and neural emulators such as Starfish, The Payne, and later neural
architectures replaced that grid with smooth interpolation of physical spectra
\citep{Czekala2015,Ting2019,Rozanski2025a}. We introduced The Payne to make
simultaneous full-spectrum inference from self-consistent ab initio models
computationally practical. A neural emulator was the means of accelerating
that calculation, not its scientific objective. The name nevertheless became
shorthand for the emulator rather than the physical fitting framework.
\end{EN}

\paynezero\ 回归这一最初目标，且不使用学习到的“标签到流量”插值器：“Zero”表示零层光谱模拟。每条试验光谱都针对所请求的标签与混合物直接合成，被接受的解用物理收敛大气检验。全光谱拟合因此可以保留所有标签的耦合响应，而无需构建高维光谱训练网格。

\begin{EN}
\paynezero\ returns to that original goal without a learned label-to-flux
interpolator. "Zero" denotes zero spectral-emulation layers. Every trial
spectrum is calculated by direct synthesis for the requested labels and
mixture, and the accepted solution is checked with a physically converged
atmosphere. Full-spectrum fitting can therefore retain the coupled response of
all labels without constructing a high-dimensional spectral training grid.
\end{EN}

\subsection{从学习流量关系到物理定标 From learned flux relations to physical calibration}

数据驱动的“标签到流量”关系从另一方向处理观测—模型差距。例如 The Cannon 直接从光谱与参考标签学习，不求值显式的物理前向模型。这条经验路线能复现合成库中缺失的可重复结构，其隐患是归因：训练样本中相关的标签可能产生有用的光谱预测器，却未必识别出某个特征背后的物理过程。

\begin{EN}
Data-driven label-to-flux relations address the observation--model gap from the
other direction. The Cannon, for example, learns directly from spectra and
reference labels rather than evaluating an explicit physical forward model.
This empirical route can reproduce repeatable structure missing from a
synthetic library. Its caveat is attribution. Correlated labels
in the training sample can produce a useful spectral predictor without
necessarily identifying the physical process responsible for a feature.
\end{EN}

当目标是改进前向模型时，直接从头算拟合检验的是被提议的物理原因。结构化残差可来自谱线数据、缺失不透明度以及更广的合成—观测失配 \citep{Czekala2015,OBriain2021}；大气、连续谱、地球大气与数据约减误差也能产生类似残差，但它们进入计算的位置不同。可微分的物理模型暴露这些输入，并通过不透明度与辐射转移检验被提议的原因。

\begin{EN}
When the goal is to improve the forward model, direct ab initio fitting tests
the proposed physical cause. Structured residuals can arise from line data,
missing opacity, and broader synthetic--observed mismatch
\citep{Czekala2015,OBriain2021}. Atmosphere, continuum, telluric, and reduction
errors can produce related residuals, but they enter the calculation in
different places. A
differentiable physical model exposes these inputs and tests a proposed cause
through opacity and radiative transfer.
\end{EN}

谱线表定标演示了这条路线：我们一次优化超过 $10^5$ 个耦合的振子强度与阻尼修正（包括混合跃迁），在单张 H100 上约一分钟得到共享太阳—大角星解。这些值是有效的标准星定标而非实验室测量；它们在冻结的 APOGEE 对照集上的表现检验了迁移效果，而未用 ASPCAP 标签来决定修正。

\begin{EN}
The line-list calibration demonstrates this route. We optimize more than
$10^5$ coupled oscillator-strength and damping corrections at once, including
blended transitions, and obtain the shared Sun--Arcturus solution in about one
minute on one H100. These values are effective standard-star calibrations rather
than laboratory measurements. Their performance on a frozen APOGEE control set
tests transfer without using ASPCAP labels to determine the corrections.
\end{EN}

生产级定标可以使用多个基准标准星 \citep{Heiter2015,Jofre2015,Smith2021}。我们建议联合求解它们的共享谱线参数，同时让每颗星的标签、连续谱与仪器响应自由变化。标准星之间的合适权重不由优化器决定：太阳与大角星探测不同的压强与分子区间，本文的等权只演示容量，而非定义一个普适定标目录。生产分析应针对被研究的恒星族群选择其标准星、权重、实验室先验与留出测试。

\begin{EN}
A production calibration could use several benchmark standards
\citep{Heiter2015,Jofre2015,Smith2021}. We propose solving their shared line
parameters while allowing each star's labels, continuum, and instrument
response to vary. The appropriate balance among standards is not determined by the
optimizer. The Sun and Arcturus probe different pressure and molecular regimes,
so the equal weighting used here demonstrates capacity rather than defining a
universal calibrated catalog. A production analysis should choose its standards,
weights, laboratory priors, and held-out tests for the stellar population being
analyzed.
\end{EN}

大幅度的不透明度修订随后会触发新的大气计算与新一轮谱线数据更新。这一外层迭代保留了物理耦合，并能暴露谱线定标无法修复的差异——包括大气迭代本身的失败。

\begin{EN}
Large opacity revisions would then trigger a new atmosphere
calculation and another line-data update. This outer iteration preserves
physical coupling and reveals discrepancies that line calibration cannot
repair, including failures of the atmosphere iteration itself.
\end{EN}

\subsection{大气边界与约化求解器 Atmosphere boundaries and reduced solvers}\label{sec:discussion_fixed_point}

一致性检验覆盖热矮星、太阳、红巨星与 K 矮星区间，巡天拟合则限于巨星。这些结果不意味着同一数值策略在每个恒星区间都保持高效。两个边缘案例指出了可能需要不同表示或求解器的位置。

\begin{EN}
The parity tests span hot-dwarf, solar, red-giant, and K-dwarf regimes, while
the survey fit is restricted to giants. These results do not imply that one
numerical strategy will remain efficient across every stellar regime. Two edge
cases identify where a different representation or solver may be needed.
\end{EN}

第一个是大气初始化器的采样维度。直接丰度模型把六条相关深度廓线压缩成紧凑 PCA 表示，但其输入仍是 3 个恒星参数加 81 个独立丰度，84 维输入相对其维度而言采样稀疏——同样的采样问题也限制高维模型网格。相关的大气响应仍让初始化器可用。未来的初始化器可以使用更贴近电子供给、分子平衡与谱线覆盖的坐标，或用化学感知的编码器学习这样的分组。初始化器只需提出初始状态，因为物理求解器仍在所请求的逐元素丰度下确立被接受的大气。

\begin{EN}
The first is the sampling dimension of the atmosphere initializer. The
direct-abundance model compresses six correlated depth profiles into a compact
PCA representation, but its input still spans three stellar parameters
together with 81 independently varied abundances. The resulting 84-coordinate
input is sparse relative to its dimension. The same sampling problem limits
high-dimensional model grids. Correlated atmospheric
responses nevertheless make the initializer useful. Future
initializers could use coordinates tied more closely to electron donation,
molecular balance, and line blanketing, or learn such groupings with a
chemistry-aware encoder. The initializer need only propose a starting state
because the physical solver still establishes the accepted atmosphere at the
requested individual abundances.
\end{EN}

第二个是冷分子大气的数值区间变化。一些更冷的高重力对照即使从邻近初始廓线出发也发展出增长的流量与加热残差，而另一些邻近起点却收敛——收敛性在该区间对初始结构变得更敏感。原 Kurucz 计算与 \paynezero\ 在其中一些对照上同样不通过收敛判据，这说明问题出在物理迭代本身，而非模拟器特有的误差。

\begin{EN}
The second is a change of numerical regime in cold molecular atmospheres. Some
cooler high-gravity controls develop growing flux and heating residuals even
from a nearby initial profile, while other nearby starts converge. Convergence
therefore becomes more sensitive to the initial structure across this regime.
Both the original Kurucz calculation and \paynezero\ fail the same convergence
criterion in some of these controls, which points to the physical iteration
rather than an emulator-specific error.
\end{EN}

\begin{physbox}{冷星大气为什么难收敛}
温度降到约 4000 K 以下，大气物理发生质变：分子（TiO、H$_2$O、CO 等）大量形成并主导不透明度；分子平衡对温度极其敏感（指数依赖）；对流层可向上延伸到直接形成光谱的层次。于是一个小的温度修正会同时重组分子平衡、不透明度、物态方程和绝热梯度，并移动对流边界——下一趟修正可能不减反增。用不动点语言说：该区间雅可比谱半径越过 1，经典逐次迭代失去局部稳定性。这是 Kurucz 类型求解器的已知固有边界，本文的重要贡献是诚实地标出了它，而不是声称学习型初值能绕过它。
\end{physbox}

该区间的小温度变化会一起重组分子平衡 \citep{BarklemCollet2016}、不透明度、物态方程与绝热梯度，对流也可能触及决定出射流量的层次 \citep{Gustafsson2008}。一趟迭代随后可能把罗斯兰不透明度、柱质量映射与对流边界移动得足以让下一步修正增大而非缩小——通常的不动点迭代失去了局部稳定性。

\begin{EN}
Small temperature changes in this regime reorganize molecular equilibrium
\citep{BarklemCollet2016}, opacity, the equation of state, and the adiabatic
gradient together. Convection can also reach the layers that set the emergent
flux \citep{Gustafsson2008}. One iteration
can then shift the Rosseland opacity, column-mass mapping, and convective
boundary enough that the next correction grows rather than shrinks. The usual
fixed-point iteration has lost local stability.
\end{EN}

在约化的平面平行表述中，我们选择温度 $T(\tau_{\rm R})$ 与单调柱质量 $m(\tau_{\rm R})$ 作为两条被迭代廓线。固定 $T_{\rm eff}$、$\log g$、$\xi$、丰度与罗斯兰网格后，流体静力学平衡与物态方程即可恢复压强、密度、电子密度与布居 \citep{HubenyMihalas2014}；不透明度设定柱质量与罗斯兰深度的关系，辐射转移设定加速度。约化求解器将强制两条耦合的根条件——能量平衡，以及柱质量与罗斯兰光学深度的一致性：
\begin{equation}
\begin{aligned}
\mathcal{R}_{T}[T,m]
&=F_{\rm rad}[T,m]+F_{\rm conv}[T,m]
-\sigma_{\rm SB}T_{\rm eff}^{4}=0,\\
\mathcal{R}_{m}[T,m]
&=\frac{\mathrm{d}\tau_{\rm R}}{\mathrm{d}m}-\kappa_{\rm R}[T,m]=0 .
\end{aligned}
\end{equation}
\noindent 这些残差提示在 $T$ 与 $m$ 中求解，而非为全部六条存档廓线做初始化。重算其余四个场并未让学习型冷矮星 $T$、$m$ 廓线恢复收敛，说明不稳定并非简单因为“预测的场太多”。

\begin{EN}
For a reduced plane-parallel formulation, we choose temperature
$T(\tau_{\rm R})$ and monotonic column mass $m(\tau_{\rm R})$ as the two
iterated profiles. At fixed $T_{\rm eff}$, $\log g$, $\xi$, abundances, and
Rosseland grid, hydrostatic balance and the equation of state then recover
pressure, density, electron density, and populations
\citep{HubenyMihalas2014}. Opacity
sets the relation between column mass and Rosseland depth, while radiative
transfer sets the acceleration. A reduced solver would enforce two coupled root
conditions, energy balance and consistency between column mass and Rosseland
optical depth,
\begin{equation}
\begin{aligned}
\mathcal{R}_{T}[T,m]
&=F_{\rm rad}[T,m]+F_{\rm conv}[T,m]
-\sigma_{\rm SB}T_{\rm eff}^{4}=0,\\
\mathcal{R}_{m}[T,m]
&=\frac{\mathrm{d}\tau_{\rm R}}{\mathrm{d}m}-\kappa_{\rm R}[T,m]=0 .
\end{aligned}
\end{equation}
\noindent These residuals motivate a solver in $T$ and $m$ rather than an
initializer for all six stored profiles. Recomputing the other four fields did
not restore convergence from learned cool-dwarf $T$ and $m$ profiles, so the
instability is not simply caused by predicting too many fields.
\end{EN}

原生的双廓线求解器必须对每组试探 $T$、$m$ 重算物态方程、不透明度、转移与对流，再联立求解两条残差。在探索性计算中，邻近物理模型估计两条残差随 $T$、$m$ 的变化并给出局部修正；残差确实下降了，但所需的完整大气求值次数多于“学习型六场初值 + 物理求解器迭代”。因此我们保留后者作为默认。若双廓线求解器的响应可以在不做多次完整物理求值的情况下估计，它仍是可选的备用方案。不同恒星区间最终可能需要不同的、经过验证的修正器。何时切换区间、同时保持同一套物理验收测试，成为一个工作流问题——它要求软件边界暴露物理状态、收敛性与验证，而不是藏在一个不透明计算内部。

\begin{EN}
A native two-profile solver must recompute the equation of state, opacity,
transfer, and convection for each trial $T$ and $m$, then solve both residuals
together. In exploratory calculations, nearby physical models estimated how
both residuals change with $T$ and $m$ and supplied local corrections. The
residuals fell, but the calculation required more complete atmosphere
evaluations than the learned six-field start followed by the physical solver iteration.
We therefore retain the latter as the default. A two-profile solver remains a
possible fallback if its response can be estimated without many full physical
evaluations. Different stellar regimes may ultimately require different
validated correctors. Choosing when to switch regimes, while preserving the
same physical acceptance tests, becomes a workflow problem. It motivates
software boundaries that expose physical state, convergence, and validation
rather than hiding them inside one opaque calculation.
\end{EN}

\subsection{作为科学接口的可读软件 Readable software as a scientific interface}

在数值区间之间安全切换，要求物理状态、残差与验收测试保持可见。原 Kurucz 工作流程把大量物理编码在由中转文件连接的紧凑 Fortran 程序中 \citep{Kurucz2005}。重实现过程中我们发现，简写的标识符、超长的例程与隐式共享状态让单个约定难以隔离：关于谱线放置、不透明度累积或状态传递的假设即使改变最终光谱，也难以审计与扩展。

\begin{EN}
Switching safely between numerical regimes requires the physical state,
residuals, and acceptance tests to remain visible. The original Kurucz workflow
encodes extensive physics in compact Fortran programs connected by staged files
\citep{Kurucz2005}. During reimplementation, we found that terse
identifiers, long routines, and implicit shared state made individual
conventions difficult to isolate. Assumptions about line placement, opacity
accumulation, or state transfer can therefore be difficult to audit and extend
even when they change the final spectrum.
\end{EN}

\paynezero\ 用具名物理数组、显式单位与阶段局部接口替代了简写标识符与隐式交接。这些选择把数值约定暴露给审计、复现、教学与扩展。

\begin{EN}
\paynezero\ replaces terse identifiers and implicit handoffs with named physical
arrays, explicit units, and stage-local interfaces. These choices expose
numerical conventions for audit, reproduction, teaching, and extension.
\end{EN}

同样的显式接口使 \paynezero\ 适合未来的智能体（agentic）工作流。Model Context Protocol 工具或任务专属技能可以把大气求解、光谱合成、谱线定标与验证暴露为带记录输入、输出与溯源的类型化操作。\footnote{见 Model Context Protocol 规范：\url{https://modelcontextprotocol.io/specification/2025-06-18/server/index}。}智能体可以检查收敛与流量残差，识别出冷的高重力模型进入了不同的数值区间，选择预验证的求解器或参数化，并启动相关诊断测试。智能体负责编排这些路径，而非为自己的结果背书；确定性的物理、数值与溯源门槛仍负责决定接受与否。

\begin{EN}
The same explicit interfaces make \paynezero\ suitable for future agentic
workflows. Model Context Protocol tools or task-specific skills could expose
atmosphere solving, spectrum synthesis, line calibration, and validation as
typed operations with recorded inputs, outputs, and provenance.
An agent could
inspect convergence and flux residuals, recognize that a cool high-gravity
model has entered a different numerical regime, select a prevalidated solver or
parameterization, and launch the relevant diagnostic tests. The agent would
orchestrate these routes rather than certify its own result. Deterministic
physical, numerical, and provenance gates would still decide acceptance.
\end{EN}

